% ===========================================================================
\section{极大理想与构造域}\label{sec:maximal}

\subsection{PID 中的极大理想}

回忆上节课的判别法（含幺交换环中）：
\[
  M\text{ 是极大理想} \iff R/M\text{ 是域}; \qquad
  P\text{ 是素理想} \iff R/P\text{ 是整环}.
\]
（域是整环，故极大理想必是素理想。）这给出构造域的方法：找一个整环，
取极大理想 $\mathfrak{m}$，则 $R/\mathfrak{m}$ 是域。在 PID 中判断极大
理想非常容易：

\begin{theorem}[PID 中极大理想的判别]
设 $R$ 为 PID，$a\in R$ 非零且非单位，则
\[
  a\text{ 不可约} \iff (a)\text{ 是极大理想}.
\]
\end{theorem}

\begin{proof}[双向都用逆否]
\leavevmode
\begin{itemize}[leftmargin=1.3em,itemsep=2pt]
  \item 若 $(a)$ 不极大，则存在 $J\supsetneq(a)$；$R$ 是 PID，$J = (b)$。
        由 $(a)\subsetneq(b)$ 知 $b\mid a$ 且不相伴，即 $b$ 是 $a$ 的
        真因子——$a$ 可约；
  \item 若 $a$ 可约，取真因子 $b$，则 $(b)\supsetneq(a)$，$(a)$ 不极大。
\end{itemize}
\end{proof}

\begin{corollary}[模掉不可约多项式得到域]
设 $F$ 为域，$f\in F[x]$ 不可约，则
\[
  F[x]/(f)\ \text{是一个域}.
\]
\end{corollary}

\begin{remark}[统一的观点]
之前见过的域也可以统一看成这种「模掉一个理想」的形式：例如
$\Q(\sqrt{2})\cong\Q[x]/(x^2-2)$——模掉理想相当于「给 $x$ 加一个
限制」：理想中的元素都被视为 $0$，即 $x^2 - 2$ 被看作 $0$，于是 $x$
就变成了 $\sqrt{2}$。
\end{remark}

\subsection{分式域的泛性质（复习与补充）}

\block{复习（分式域）}
设 $R$ 为整环，所有「分数」 $\frac{a}{b}$（$b\ne 0$）模去等价关系
$\frac{a}{b}\sim\frac{c}{d} \iff ad = bc$，所得集合构成域 $\Frac(R)$。

\begin{theorem}[分式域的泛性质 / 最小性]
$\Frac(R)$ 是把 $R$ 扩成域的\keyword{最小}办法：设
$\varphi\colon R\to R'$ 是单同态，且对任意 $x\ne 0$ 都有
$\varphi(x)\in(R')^\times$（$R'$ 是域时此条件自动成立），则存在
\keyword{唯一}同态 $\tilde\varphi\colon\Frac(R)\to R'$ 使
\[
  \tilde\varphi\circ\iota = \varphi, \qquad
  \iota\colon R\hookrightarrow\Frac(R),\quad x\mapsto\tfrac{x}{1}.
\]
\end{theorem}

\begin{proof}[构造]
$\tilde\varphi$ 的定义别无选择：
\[
  \tilde\varphi\!\left(\frac{a}{b}\right) = \varphi(a)\,\varphi(b)^{-1},
\]
这成立是因为 $\varphi(b)$ 在 $R'$ 中可逆。再验证它确实是同态即可。
\end{proof}
